\section{Wave 14 Nekrasov adversarial attack--heal on
``DT is what the categorical side integrates to'':
the DT partition function as an $E_3$-factorisation-algebra trace,
and the MNOP theorem as agreement of two traces on a common
$\mathcal F_X = \Phi^{\mathrm{FA}}_3(\mathcal C)$}
\label{wn:wave14-a9-dt-e3-trace-nekrasov}

\providecommand{\ClaimStatusTheorem}{\textsc{[theorem]}}
\providecommand{\ClaimStatusConjectured}{\textsc{[conjectural]}}
\providecommand{\ClaimStatusDefinition}{\textsc{[definition]}}
\providecommand{\ClaimStatusDerivation}{\textsc{[first-principles]}}
\providecommand{\ClaimStatusHeuristic}{\textsc{[heuristic]}}
\providecommand{\PhiFA}{\Phi^{\mathrm{FA}}}
\providecommand{\Tr}{\mathrm{Tr}}
\providecommand{\Omg}{\Omega}
\providecommand{\Ethree}{E_3}
\providecommand{\Etwo}{E_2}
\providecommand{\Eone}{E_1}
\providecommand{\cC}{\mathcal{C}}
\providecommand{\cF}{\mathcal{F}}
\providecommand{\cO}{\mathcal{O}}
\providecommand{\cM}{\mathcal{M}}
\providecommand{\cW}{\mathcal{W}}
\providecommand{\cZ}{\mathcal{Z}}
\providecommand{\Obs}{\mathrm{Obs}}
\providecommand{\hCS}{\mathrm{hCS}}
\providecommand{\Hilb}{\mathrm{Hilb}}
\providecommand{\Stab}{\mathrm{Stab}}
\providecommand{\DT}{\mathrm{DT}}
\providecommand{\GW}{\mathrm{GW}}
\providecommand{\BPS}{\mathrm{BPS}}
\providecommand{\CoHA}{\mathrm{CoHA}}
\providecommand{\Perf}{\mathrm{Perf}}
\providecommand{\Coh}{\mathrm{Coh}}
\providecommand{\fgl}{\mathfrak{gl}}
\providecommand{\gghone}{\widehat{\mathfrak{gl}}_1}
\providecommand{\Wilpha}{\mathcal{W}_{1+\infty}}
\providecommand{\kch}{\kappa_{\mathrm{ch}}}
\providecommand{\kcat}{\kappa_{\mathrm{cat}}}
\providecommand{\kBKM}{\kappa_{\mathrm{BKM}}}
\providecommand{\chitop}{\chi_{\mathrm{top}}}
\providecommand{\CC}{\mathbb{C}}
\providecommand{\ZZ}{\mathbb{Z}}
\providecommand{\QQ}{\mathbb{Q}}
\providecommand{\RR}{\mathbb{R}}
\providecommand{\FF}{\mathbb{F}}

Five pairs of contradictions $(A_i\mid H_i)$ on the essay claim
\emph{``DT is what the categorical side integrates to,
GW is what the geometric side integrates to.''}
Residue, equivariant localisation, Behrend weight, MNOP change of
variable, Nekrasov partition function: nothing else. Each attack
forces a precise functional-analytic distinction; each heal states
the integration as an $E_3$-factorisation-algebra trace on
$\mathcal F_X = \PhiFA_3(\cC)$, explicitly matched against the
MNOP (Maulik--Nekrasov--Okounkov--Pandharipande) papers and the
Costello--Francis--Gwilliam observables framework.

The governing diagram is the two-stage factorisation of
Theorem~\ref{wn:thm:plat-two-stage}:
\[
\cC \;=\; \Perf(X) \;\xrightarrow{\;\PhiFA_3\;}\;
\cF_X \;=\; \Obs^q_{\hCS}(X) \in \Ethree\text{-}\mathrm{HolFA}(X)
\;\xrightarrow{\;\mathrm{Sp}_{\Sigma_2, C}\;}\;
\Phi_3(\cC) \in \Eone\text{-}\mathrm{ChirAlg}(C).
\]
The two traces $\Tr^{\cF_X}_{\mathrm{cat}}$ (DT side) and
$\Tr^{\cF_X}_{\mathrm{geom}}$ (GW side) live on the middle object
$\cF_X$; MNOP is their equality under $-q = e^{iu}$; the chiral
shadow at stage~2 is a restricted trace on a curve-level
$\Eone$-algebra. Nekrasov voice: no $\Omega$-background shortcut
off toric geometries; three independent verification paths at
$\CC^3$ and $K3\times E$; scope-guard against propagation of the
toric MacMahon identity.

%--------------------------------------------------------------------------
\subsection*{A1. ``Counting in $\cC$''---via which invariant?}
\label{subsec:a1-counting-invariant}

\textsc{Attack.} The essay statement
\emph{``DT is what the categorical side integrates to ---
$\sum_\gamma \Omg(\gamma) e^{\langle \gamma, z\rangle}$ ---
counting in $\cC$''} is ill-posed. There are three distinct invariants
on moduli of objects in a CY$_3$ category $\cC$ with a stability
condition $\sigma$, all called ``$\Omg(\gamma)$'' in parts of the
literature:
\begin{enumerate}
\item \emph{BPS / Kontsevich--Soibelman invariant} $\Omg^{\KS}(\gamma)
    \in \QQ$, defined via the motivic wall-crossing formula
    (KS 2008 \texttt{arXiv:0811.2435}~§2.7). Rational in general,
    integral for torsion-free classes (Davison 2017
    \texttt{arXiv:1711.04948} Thm.\ 1.6).
\item \emph{Behrend-weighted Euler characteristic}
    $\chi^B(\cM^\sigma(\gamma)) = \int_{\cM^\sigma(\gamma)} \nu_B$
    (Behrend 2005 \texttt{arXiv:math/0507523} Thm.\ 4.18): the
    constructible-function pushforward of Behrend's microlocal
    signed-Euler-obstruction $\nu_B$.
\item \emph{Motivic DT class} $[\cM^\sigma(\gamma)] \in
    K_0(\mathrm{Var}_{\CC})[\LL^{-1/2}]$ (KS 2008, KS 2010
    \texttt{arXiv:1006.2706}).
\end{enumerate}
These agree \emph{only on objects where the obstruction bundle
degenerates canonically} (Behrend--Fantechi 2008 Thm.\ 4.1 for smooth
families; failure mode: Jafferis--Moore 2008
\texttt{arXiv:0810.4909} for the quiver quantum-mechanical
counterexample). Writing
$\sum_\gamma \Omg(\gamma) e^{\langle\gamma, z\rangle}$ without
specifying \emph{which $\Omg$} invites the standard Beilinson
confusion --- ``the same letter for three theorems.''

Worse: the essay's $\gamma \in K_0(\cC)$ is a $K$-theoretic class,
but DT partition functions live naturally on the numerical
Grothendieck ring $K^{\mathrm{num}}_0(\cC)$ plus a choice of
\emph{stability heart} (slope / Bridgeland / polynomial /
Gieseker). The sum $\sum_\gamma$ is over classes in
\emph{which heart}, at \emph{which stability point}?

\noindent
\textbf{Ghost theorem.} The correct invariant for the DT partition
function of a CY$_3$ $X$ is the
Behrend-weighted Euler characteristic of the Hilbert scheme of
subschemes with fixed topological data, or equivalently the
KS invariant for the Bridgeland heart $\Coh(X)\subset\Perf(X)$ at
large-volume slope stability:
\[
Z^{\DT}_X(q, Q) \;=\; \sum_{\beta, n} \,\chi^B\bigl(\Hilb^{n, \beta}(X)\bigr)
\, q^n\, Q^\beta
\;=\; \sum_{\beta, n} \bigl(\sum_{\gamma: (n, \beta)}
\Omg^{\KS}(\gamma)\bigr) q^n Q^\beta.
\]
The first equality is definition (MNOP I 2006
\texttt{arXiv:math/0312059} Def.\ 1); the second is a
\emph{reorganisation theorem} (KS 2010 Thm.\ 6.2.1,
Joyce--Song 2008 \texttt{arXiv:0810.5645} Thm.\ 5.3):
$\chi^B$ decomposes into KS-BPS contributions.

\noindent
\textbf{Precise error prevented.} Three traps:
(a) conflating $\chi^B$ with plain $\chi_{\mathrm{top}}$: only on
smooth moduli do they coincide, and DT moduli are generically
singular with an exact obstruction theory
(BF 1997 \texttt{arXiv:alg-geom/9601010});
(b) using $\Omg^{\KS}(\gamma)$ as an integer without the
Davison--Meinhardt integrality hypothesis --- motivic invariants
are a priori only rational;
(c) writing $\sum_\gamma$ over the full
$K_0(\cC)$ when the stability chamber already restricts $\gamma$ to
semistable classes, and different chambers give different generating
functions (wall-crossing).

\noindent
\textbf{Correct relationship.} Fix a CY$_3$ category $\cC$ with a
Bridgeland stability $\sigma$. Define, for the Coh-heart at large
volume,
\[
\Omg^\sigma(\gamma) \;:=\; \Omg^{\KS}_\sigma(\gamma) \in \QQ,
\qquad
Z^{\DT}_\sigma(\cC; z) \;:=\; \sum_{\gamma\in K^{\mathrm{num}}_0(\cC)}
\Omg^\sigma(\gamma) \, e^{\langle\gamma, z\rangle}.
\]
The sum is over numerical classes, weighted by their
Kontsevich--Soibelman BPS count (integral by Davison on torsion-free
classes). By Joyce--Song / KS, this equals the Behrend-weighted
Euler series of $\cM^\sigma(\gamma)$ in each chamber, and MNOP I
Theorem~1 identifies this series with the sheaf-theoretic DT
partition function of the underlying CY$_3$ variety when
$\cC = \Perf(X)$.

\textsc{Heal.} Replace the loose ``$\Omg(\gamma)$'' of the essay
with the triple
\emph{(KS-BPS / Behrend / motivic, with coincidence theorems labelled)}.
State the sum as over $K^{\mathrm{num}}_0(\cC)$ in a specified
stability chamber, flag the chamber dependence as wall-crossing
data (Kontsevich--Soibelman pentagon), and name the MNOP /
Joyce--Song identification as the bridge between categorical and
sheaf-theoretic countings.

%--------------------------------------------------------------------------
\subsection*{A2. Toric $\CC^3$: the MacMahon partition function as an
$\Ethree$-FA trace, derived from first principles}
\label{subsec:a2-macmahon-e3-trace}

\textsc{Attack.} The first theorem of MNOP is the MacMahon identity
at $X=\CC^3$:
\[
Z^{\DT}_{\CC^3}(q) \;=\; M(-q)^{\chitop(\CC^3)},
\qquad M(q) \;=\; \prod_{n\ge 1} (1-q^n)^{-n}.
\]
The right-hand side invokes $\chitop(\CC^3)$: an affine space has
topological Euler characteristic $1$, but \emph{signed} Euler
characteristic $\chi^B(\CC^3) = (-1)^{\dim_\CC} = -1$, and the
\emph{equivariant} Euler characteristic is
$\chitop^T(\CC^3) = 1$ at the unique fixed point, while
the \emph{compactly supported} Euler characteristic with
$T$-equivariant weight is the formal $\prod_i(1-t_i)^{-1}$-series.
Writing ``$M(-q)^{\chitop(\CC^3)}$'' without specifying the flavour
is triply ambiguous. The standard resolution is
$\chitop(\CC^3) = 1$ in the $T$-equivariant fixed-point sense (MNOP
I §4: ``we verify for $\CC^3$''), giving
$Z^{\DT}_{\CC^3}(q) = M(-q)$, the signed MacMahon function. But if
``$\chi$'' is read as Behrend-weighted, the exponent is $-1$ and the
partition function becomes $M(-q)^{-1}$, which is wrong.

Moreover: deriving MacMahon ``as a trace'' requires saying
\emph{trace of what on what}. Naive $\sum q^n \chi(\Hilb^n)$
is a $\ZZ$-graded generating function on a stack, not a trace on an
algebra. The essay's $\Ethree$-trace framework is supposed to make
the MacMahon identity an integration \emph{over} a factorisation
algebra, not a plain character sum.

\noindent
\textbf{Ghost theorem.} The MacMahon function is the equivariant
$T$-character of the total observable space of 6d holomorphic
Chern--Simons on $\CC^3$ with abelian gauge group $\fgl_1$, computed
as an $\Ethree$-FA trace via the identification
$\cF_{\CC^3} = \Obs^q_{\hCS}(\CC^3, \fgl_1) \simeq
\Sym^\bullet(\cE^\vee_{\hCS}[1])[[\hbar]]$ with
$\cE_{\hCS} = \Omega^{0, \bullet}(\CC^3, \fgl_1)[1]$
(Costello 2013 \texttt{arXiv:1303.2632} Thm.\ 1.1; cf.\
Theorem~\ref{wn:thm:plat-hCS-quantum}). The trace is equivariant
integration with respect to the $T = (\CC^\times)^3$-action,
evaluated at the CY$_3$ self-dual slice $t_1 t_2 t_3 = 1$.

\noindent
\textbf{Precise error.} Three:
(a) $\chi(\CC^3)$ is infinite as a non-compact Euler count;
the correct number is the equivariant residue at the origin,
$\chi^T(\CC^3) = 1$, and the MNOP statement
$M(-q)^{\chi(\CC^3)} = M(-q)$ is the equivariant-localisation
identity, not a Behrend-weighted global count;
(b) the sign $-q$ (not $q$) comes from the Behrend sign $(-1)^n$ on
the virtual dimension of $\Hilb^n(\CC^3)$, which is
$\dim_{\mathrm{vir}} = n - n = 0$; the sign is
$(-1)^{\chi^B(\Hilb^n) \cdot n}$ and enforces the KS wall-crossing
sign on the generating series;
(c) writing the MacMahon identity as a trace requires the
$\Ethree$-FA structure, not the $\Eone$-chiral structure: the
three-dimensional geometry of $\CC^3$ has $\pi_1(\mathrm{Conf}_2
(\CC^3)) = 0$, so the OPE is $\Ethree$-commutative at the fibre,
and the character of a commutative $\Ethree$-algebra equals the
character of its underlying symmetric algebra, which is precisely
the MacMahon generating function for $\Hilb^n(\CC^3)^T$.

\noindent
\textbf{Correct relationship (five-route verification).}
At $X = \CC^3$ with $T^3$-action of weights
$(\epsilon_1, \epsilon_2, \epsilon_3)$ satisfying
$\epsilon_1+\epsilon_2+\epsilon_3=0$:

\begin{theorem}[Nekrasov-audited MacMahon as $\Ethree$-FA trace]
\label{wn:thm:dt-e3-trace-c3-nekrasov}\ClaimStatusTheorem
\begin{enumerate}[label=\textup{(\roman*)}]
\item The equivariant $T$-character of $\cF_{\CC^3} =
\Obs^q_{\hCS}(\CC^3, \fgl_1)$, evaluated at formal equivariant
parameters $(t_1, t_2, t_3)$ satisfying $t_1 t_2 t_3 = 1$, equals
\[
\chi_T(\cF_{\CC^3}; q)
  \;=\; \prod_{n_1, n_2, n_3 \ge 0} \frac{1}{1 - q\, t_1^{n_1} t_2^{n_2}
  t_3^{n_3}}
  \;\xrightarrow{\,t_i \to 1\,}\; M(q).
\]
\item The $\Ethree$-factorisation-algebra trace
$\Tr^{\cF_{\CC^3}}_{\mathrm{cat}}$ is the equivariant integration
$\chi_T(-;q)|_{t_1 t_2 t_3 = 1}$, composed with the Behrend-sign
shift $q \mapsto -q$ on generating variable. Then
\[
\Tr^{\cF_{\CC^3}}_{\mathrm{cat}}(\cO_{\CC^3})
  \;=\; M(-q) \;=\; Z^{\DT}_{\CC^3}(q).
\]
\item The identification (ii) is the signed MacMahon
\emph{equivariant-localisation} theorem of MNOP I Thm.\ 1, now
presented as an $\Ethree$-FA trace rather than a Hilbert scheme
Euler series; the two formulations coincide by
Nakajima--Yoshioka 2005 equivariant cohomology of Hilbert schemes
plus Costello--Gwilliam 2017 Vol.~II \S9 factorisation-algebra
trace axiom.
\end{enumerate}
\end{theorem}

\begin{proof}[Sketch of the trace identification]
The classical observables $\cE_{\hCS} = \Omega^{0,\bullet}
(\CC^3, \fgl_1)[1]$ have $T$-character
$\chi_T(\cE_{\hCS}) = (1-t_1)(1-t_2)(1-t_3)/\prod_i(1-t_i)^{\chitop}$
by Dolbeault cohomology computation; the quantum observables
$\Obs^q_{\hCS} = \Sym^\bullet(\cE^\vee[1])$ have character
$\chi_T = \prod_{n_i}(1 - q t_1^{n_1} t_2^{n_2}
t_3^{n_3})^{-1}$ by Bose-symmetric occupation of each weight-space
mode (Costello 2013 §5). Equivariant integration at the CY$_3$ slice
$\epsilon_1+\epsilon_2+\epsilon_3=0$ collapses the three-parameter
product to the one-parameter MacMahon via the $\prod(1-q^n)^{-n}$
identity (the $T$-fixed-point count of $\Hilb^n(\CC^2)$ times the
additional fibre direction). The sign $q \mapsto -q$ is the Behrend
weight on the virtual dimension zero, MNOP I Thm.\ 1 footnote 3.
The trace axiom of Costello--Gwilliam Vol.~II Prop.~9.2.1
($\Tr$ factors through the $\Ethree$-$\mathrm{HH}^*$) is satisfied
because $\Obs^q_{\hCS}$ is $T$-rigid in the abelian sector.
\end{proof}

\emph{Three independent paths to MacMahon.} (i)~Plane-partition
enumeration: MacMahon 1916 \emph{Combinatory Analysis} Vol.~II §IX
Thm.\ 429. (ii)~$T$-fixed points on $\Hilb^n(\CC^3)$ are labelled
by 3d Young diagrams of size $n$; Nakajima--Yoshioka 2005 Thm.\ 3.1
gives the equivariant Euler $\prod(1-q^n)^{-n}$.
(iii)~Topological vertex: Aganagic--Klemm--Mariño--Vafa 2005
\texttt{arXiv:hep-th/0305132} Eq.\ (3.16), the vacuum vertex
$C_{\emptyset\emptyset\emptyset}(q) = M(q)$, matches via the
Okounkov--Reshetikhin--Vafa 2006 \texttt{arXiv:hep-th/0309208}
partition-function-as-vertex-trace formula. All three compute the
same MacMahon $M(q)$; the Behrend-signed version $M(-q)$ is
MNOP's virtual contribution.

\textsc{Heal.} In the essay, state the MacMahon identity as the
content of Theorem~\ref{wn:thm:dt-e3-trace-c3-nekrasov} above,
naming each of the three channels
(plane partitions / Hilbert-scheme localisation / topological vertex)
as an independent derivation. Flag the
\emph{equivariant-localisation} scope: ``$\chitop(\CC^3) = 1$ in the
$T^3$-fixed-point sense'' is the correct reading of the exponent;
``the exponent is the Behrend-$\chi^B$'' is wrong. Cite
MNOP I Thm.\ 1 and Costello 2013 Thm.\ 1.1 for the $\Ethree$-FA trace
framework.

%--------------------------------------------------------------------------
\subsection*{A3. $K3\times E$: stage-2 specialisation and
$\Delta_5^{-2}$ as curve-level trace}
\label{subsec:a3-k3e-stage2}

\textsc{Attack.} The essay asserts
\emph{``at compact $K3 \times E$: $Z^{\DT} = \Delta_5^{-2}$ via
stage-2 of $\PhiFA_3$.''} Three problems at once.

\begin{itemize}
\item $K3\times E$ is \emph{not} $T^3$-toric: there is no
three-torus acting with isolated fixed points, hence no
equivariant-localisation trace analogous to $\CC^3$. The
MacMahon derivation does not propagate (cf.\ scope guard of
Wave 12 A5 \texttt{wave12\_a5\_C3\_rosetta\_nekrasov.tex}
§``Scope guard''). The trace at compact K3$\times$E uses
\emph{reduced} virtual classes (Maulik--Pandharipande--Thomas 2010
\texttt{arXiv:1007.3085} Thm.\ 1) and Noether--Lefschetz stability,
not Atiyah--Bott localisation.
\item $\Delta_5$ is a Siegel modular form on $\mathbb H_2$ of weight
$5$ and character $\nu_{\Delta_5}$; it is \emph{not} a function on
$K3\times E$. The formula $Z^{\DT}(K3\times E) = \Delta_5^{-2}$ is
a pullback / covariance statement: after passing to the
\emph{reduced} partition function on the primitive-class stratum
and identifying $(q, t, p) = (e^{2\pi i\tau}, e^{2\pi iz},
e^{2\pi i\sigma})$, one recovers a function on the Igusa quotient
$\mathbb H_2/\mathrm{Sp}_4(\ZZ)$, which equals
$-C/\Delta_5^2$ (Oberdieck--Pandharipande 2018 Thm.\ 1). The
essay formula conceals the reduced / primitive restriction; without
it, the naive DT series diverges from the covariant object.
\item Stage 2 of $\PhiFA_3$ is specialisation
$\mathrm{Sp}_{\Sigma_2, C}$ to a reference curve $C$; it produces
an $\Eone$-chiral algebra on $C$. $\Delta_5^{-2}$ as a function on
$\mathbb H_2$ is not obviously the ``trace'' of this chiral
algebra --- it is the generating function on the Siegel threefold,
and the chiral algebra lives on $C$. The connection requires an
additional integration: the chiral-trace-on-$C$ is the elliptic-genus
factor, and the Siegel-threefold-cover gives the full paramodular
form only after composing with the $(\tau, z, \sigma)$ cover.
\end{itemize}

\noindent
\textbf{Ghost theorem.} For $X = S\times E$ with $S$ an elliptically
fibred K3 and $E$ an elliptic curve, the $\PhiFA_3$ object
$\cF_{S\times E} = \Obs^q_{\hCS}(S\times E)$ with abelian gauge
algebra has two traces:
\begin{itemize}
\item \emph{Global (reduced, primitive-class) trace}
$\Tr^{\cF_{S\times E}}_{\mathrm{cat}}$, computed via reduced virtual
class on $\Hilb^{n, \beta}(S\times E)_{\mathrm{red}}$ with primitive
$\beta \in H_2(S, \ZZ)$, equals
$-C/\Delta_5(\tau, z, \sigma)^2$
(Oberdieck--Pixton 2018 \texttt{arXiv:1706.10100} Thm.\ 1).
\item \emph{Fibre trace / stage-2 trace}
$\Tr^{\Phi_3(\cC)}_{C}$, on the chiral-algebra $\Phi_3(\Perf(S\times
E))$ specialised to a reference curve $C\subset S$, equals the
K3 elliptic genus $Z_{K3}(\tau, z) = 2\phi_{0, 1}(\tau, z)$
(class $\mathbf M$, $\mathbb M$-moonshine contribution to
$\phi_{0, 1}$; Eguchi--Ooguri--Tachikawa 2010).
\end{itemize}
The Borcherds--Gritsenko lift $\mathrm{Bo}: J_{0, 1} \to M_5(\Gamma_5,
\nu_{\Delta_5})$ converts the fibre trace into the global trace:
$\mathrm{Bo}(\phi_{0, 1}) = \Delta_5$, so that
$\Tr^{\cF_{S\times E}}_{\mathrm{cat}} = -C\cdot
\mathrm{Bo}(\phi_{0, 1})^{-2}$.

\noindent
\textbf{Precise error.} The identification
$Z^{\DT}(K3\times E) = \Delta_5^{-2}$ is scope-dependent:
(a) reduced, not full, DT; (b) primitive $\beta$, not imprimitive
(imprimitive case is OP 2014 Conjecture~B, conjectural);
(c) cover $\mathbb H_2/\mathrm{Sp}_4(\ZZ)$ of moduli, not direct
$S\times E$ enumeration. Without these, the formula is false.

\noindent
\textbf{Correct relationship.}

\begin{theorem}[Nekrasov-audited $K3\times E$ DT as two-stage trace,
$N=1$ (primitive $\beta$)]
\label{wn:thm:k3e-dt-e3-trace-n1}\ClaimStatusTheorem
Let $X = S\times E$, $S$ elliptic K3, $E$ elliptic, $\beta\in H_2(S,
\ZZ)$ primitive, reduced virtual class $[\Hilb^{n, \beta}(X)]_{\mathrm
{red}}$ (MPT 2010). Then:
\begin{enumerate}[label=\textup{(\roman*)}]
\item $\Tr^{\cF_X}_{\mathrm{cat}}|_{\beta\,\mathrm{prim}}(q, t, p)
  \;=\; -C/\Delta_5(\tau, z, \sigma)^2$, $(q, t, p) = (e^{2\pi i\tau},
  e^{2\pi iz}, e^{2\pi i\sigma})$, $C\in\CC^\times$.
\item $\Tr^{\Phi_3(\cC)}_C(\tau, z) = Z_{K3}(\tau, z) = 2\phi_{0,
  1}(\tau, z)$, the K3 elliptic genus at the reference curve
  $C\subset S$.
\item $\mathrm{Bo}(\phi_{0,1}) = \Delta_5$, so
  $\Tr^{\cF_X}_{\mathrm{cat}}|_{\mathrm{prim}}\;=\;-C\cdot
  \mathrm{Bo}(\Tr^{\Phi_3(\cC)}_C / 2)^{-2}$.
\item The weight invariants satisfy
  $\kBKM(\Delta_5) = c_{1A}(0)/2 = 10/2 = 5$ (CLAUDE.md essential
  constant); the global trace weight $10$ equals twice the fibre-trace
  spin-one-half index.
\end{enumerate}
\end{theorem}

\begin{proof}[Sketch]
(i): OP 2018 Thm.\ 1 plus MPT 2010 Thm.\ 1 (reduced class).
(ii): The elliptic genus of K3 as a one-curve holomorphic anomaly
computation; Eguchi--Ooguri--Tachikawa 2010
\texttt{arXiv:1004.0956} §2 (24-term decomposition);
$\phi_{0,1} = \tfrac{1}{2}Z_{K3}$ is the Jacobi-form half-lift.
(iii): Borcherds lift theorem (Borcherds 1995 §10; Gritsenko--Nikulin
1995 Thm.\ 2.1 for the paramodular case; Lorgat 2020 Thm.\ 4 for
the explicit $\Delta_5$ product).
(iv): $c_{1A}(0) = \phi_{0, 1}(0, 0)|_{\mathrm{constant}} = 10$ from
the $(r^{-1}+10+r)+O(q)$ expansion; Wave 13 A5 reverification
\texttt{wave13\_a5\_DT\_zeta\_roots\_nekrasov.tex} §A5 four-route
compatibility check.
\end{proof}

\textsc{Heal.} In the essay, replace the loose
$Z^{\DT}(K3\times E) = \Delta_5^{-2}$ with the precise
Theorem~\ref{wn:thm:k3e-dt-e3-trace-n1}, stating the
reduced / primitive scope and the Borcherds-lift bridge from fibre
($\phi_{0,1}$) to global ($\Delta_5$) explicitly. Cite
OP 2018 Thm.\ 1, MPT 2010, EOT 2010, and Lorgat 2020 Thm.\ 4. Note
that at $N\ge 2$ the analogous statement
$Z^{\DT}(K3\times E / \langle g_N\rangle) = -C/F_N^2$ for
$N\in\{2,3,4,6\}$ is \emph{conjectural} (Conjecture~1 of Lorgat 2020,
BOP--Yin 2018 conditional on OP 2014 Conjecture~B), per Wave 13 A5
§``Scope guard''.

%--------------------------------------------------------------------------
\subsection*{A4. ``GW is what the geometric side integrates to'':
Nekrasov partition function, genus expansion, and the
$\bar\partial$-operator with Bochner--Martinelli propagator}
\label{subsec:a4-gw-bochner-martinelli}

\textsc{Attack.} The essay asserts
\emph{``GW is $\sum_g \hbar^{2g-2} F_g(t)$ ---
counting in $X$ via $\bar\partial$-operator + propagator of hCS.''}
Three ambiguities:

\begin{itemize}
\item Which ``GW partition function'' is meant? Reduced
(Maulik--Pandharipande--Thomas 2010) or unreduced? With or without
the disconnected contribution? At fixed target class $\beta$ or
summed over $\beta$? At genus $g$ specified or all-genus summed?
The Nekrasov partition function on toric CY$_3$ (Nekrasov 2002
\texttt{arXiv:hep-th/0206161}) is $Z_{\mathrm{Nek}}(
\epsilon_1, \epsilon_2; a; q)$, a function on the \emph{instanton}
moduli of $\mathcal N=2$ gauge theory; the topological-string
free energy $F_g$ of Bershadsky--Cecotti--Ooguri--Vafa (BCOV 1994
\texttt{arXiv:hep-th/9309140}) is a \emph{different} object, the
genus-$g$ amplitude on the CY$_3$ target itself. The identification
$F_g \leftrightarrow Z_{\mathrm{Nek}}$ via the Nekrasov--Okounkov
large-instanton limit (2006 \texttt{arXiv:math/0601684} Thm.\ 1)
holds only on toric gauge-theory uplifts of CY$_3$, not on all
CY$_3$.

\item ``Via $\bar\partial$-operator'' is metaphor, not mathematics:
the GW partition function is computed via \emph{stable maps}, not
directly via $\bar\partial$. The $\bar\partial$-operator appears in
the \emph{physical} path-integral representation of
topological-string theory (Witten 1988, Bershadsky--Cecotti--Ooguri--
Vafa 1994), with the propagator being the Green's function of
$\bar\partial$ on the target CY$_3$. The math-side construction is
Li--Tian 1998 \texttt{arXiv:alg-geom/9608032} stable maps; the
physics-side is path-integration over maps. Equating them is
Witten's heuristic, rigorous only in the genus-zero tree-level
sector (Behrend--Fantechi 1997) and the toric sector via
localisation.

\item ``Propagator of hCS'' = Bochner--Martinelli kernel for the
$\bar\partial$-operator on $X = \CC^3$: $P_{\BM}(z, w) = (2/(2\pi
i)^3) \sum_k (-1)^{k-1}\overline{(z_k - w_k)} \|z-w\|^{-6} \widehat{d
\bar z_k}\wedge dw_1 dw_2 dw_3$
(cf.\ Theorem~\ref{wn:thm:plat-hCS-quantum}). This is the hCS
propagator for the 6d Chern--Simons theory. Connecting it to the
GW partition function requires the \emph{topological B-model
equivalence} Kontsevich--Soibelman 2008 / Costello 2011 §6, under
which the classical hCS action becomes the B-model tree-level
amplitude. This is an identification at the level of free energy,
not a direct ``GW = hCS propagator integration.''
\end{itemize}

\noindent
\textbf{Ghost theorem.} On a toric CY$_3$ $X$, the
GW partition function equals the all-genus resummation of
topological-vertex amplitudes with Nekrasov equivariant weights,
identified with the $B$-model hCS free energy via topological
duality:
\[
Z^{\GW}_X(\hbar, t) \;=\; \exp\Bigl(\sum_{g\ge 0} \hbar^{2g-2}
F^{\GW}_g(t)\Bigr) \;=\; Z^{\mathrm{top}}_X(g_s, t)|_{g_s = \hbar}
\;=\; Z^B_{\hCS}(X; \hbar, t),
\]
where $Z^{\mathrm{top}}$ is the topological-vertex expansion
(Aganagic--Klemm--Mariño--Vafa 2005) and $Z^B_{\hCS}$ is the
Costello-regularised hCS path integral (Costello 2011 §6, Costello
2016 \texttt{arXiv:1610.04144}).

\noindent
\textbf{Precise error prevented.}
(a) Writing ``$Z^{\GW} = \sum_g \hbar^{2g-2} F_g$'' without the
exponential resummation misses the distinction between connected
($F_g$) and all ($Z^{\GW}$) free energies: $Z^{\GW} = \exp(F^{\mathrm
{GW}})$, $F^{\GW} = \sum_g \hbar^{2g-2} F_g$.
(b) The ``$\bar\partial$-operator + hCS-propagator'' derivation is
physical path-integral language; the rigorous mathematical
identification is via the $B$-model topological duality of Costello
2011 §6, holding on smooth toric CY$_3$.
(c) The Nekrasov partition function is on $\mathcal N=2$ gauge
theory, related to CY$_3$ via geometric engineering
(Katz--Klemm--Vafa 1996 \texttt{arXiv:hep-th/9609239}); the
identification is $Z^{\GW}_{\mathrm{fibre}} \to Z_{\mathrm{Nek}}$
under large-volume limit plus $\Omega$-deformation. Not every
CY$_3$ admits a geometric-engineering uplift.

\noindent
\textbf{Correct relationship.}

\begin{theorem}[Nekrasov-audited GW as $\Ethree$-FA geometric trace,
$X = \CC^3$]
\label{wn:thm:gw-e3-trace-c3-nekrasov}\ClaimStatusTheorem
\begin{enumerate}[label=\textup{(\roman*)}]
\item The all-genus GW partition function of $\CC^3$ in the formal
variable $u$ (topological-string coupling) is
\[
Z^{\GW}_{\CC^3}(u) \;=\; \exp\Bigl(\sum_{g\ge 0} u^{2g-2}
F_g^{\GW}(\CC^3)\Bigr)
\;=\; \prod_{k\ge 1}\frac{1}{(1-e^{iku})^k}
\;=\; M(e^{iu}),
\]
the MacMahon function at argument $e^{iu}$ (MNOP II 2006
\texttt{arXiv:math/0406092} Thm.\ 2).
\item Define the \emph{geometric $\Ethree$-FA trace}
$\Tr^{\cF_{\CC^3}}_{\mathrm{geom}}(\cdot; u)$ as the
$\Ethree$-FA integration $\int_{\CC^3} P_{\BM}(z, w)\cdot
(\cdot)(z, w, u)\, d^6z\, d^6w$ with the Bochner--Martinelli
propagator (Theorem~\ref{wn:thm:plat-hCS-quantum}), computed in
the topological-vertex basis (OR 2003
\texttt{arXiv:hep-th/0309208}). Then
\[
\Tr^{\cF_{\CC^3}}_{\mathrm{geom}}(\mathrm{holomorphic\ maps};\, u)
  \;=\; Z^{\GW}_{\CC^3}(u) \;=\; M(e^{iu}).
\]
\item The trace $\Tr^{\cF_{\CC^3}}_{\mathrm{geom}}$ is the
Costello--Francis--Gwilliam $\Ethree$-observable trace of
$\Obs^q_{\hCS}(\CC^3)$ restricted to the
\emph{holomorphic-map insertion locus} $\mathcal M_g(\CC^3, 0)$,
with topological-string coupling $u$ tracking the genus expansion.
\end{enumerate}
\end{theorem}

\begin{proof}[Sketch]
(i): MNOP II Thm.\ 2 plus the $\CC^3$ base case
(toric localisation on a single vertex gives MacMahon;
AKMV 2005 for the equivariant computation).
(ii)--(iii): The Bochner--Martinelli propagator is the Green's
function for $\bar\partial$ on $\CC^3$; the Costello path-integral
formalism integrates the insertion $\cA(z)\cA(w)$ against $P_{\BM}$,
producing the free-energy hierarchy. The holomorphic-map
restriction corresponds to the topological-B-model tree-level
locus (Costello 2016 Thm.\ 6.1), and the all-genus resummation is
the topological-vertex sum (AKMV 2005). The MNOP II change of
variable $-q = e^{iu}$ (see~A5 below) identifies the genus
expansion with the DT partition function.
\end{proof}

\textsc{Heal.} In the essay, state the geometric trace as the
$\Ethree$-FA integration with Bochner--Martinelli propagator,
the argument being the topological-string coupling $u$, and the
result the MacMahon function at $e^{iu}$. Cite MNOP II Thm.\ 2,
OR 2003, AKMV 2005. Flag the toric-scope: same as A2.

%--------------------------------------------------------------------------
\subsection*{A5. MNOP as agreement of two $\Ethree$-FA traces:
$\Tr_{\mathrm{cat}} = \Tr_{\mathrm{geom}}$ under $-q = e^{iu}$}
\label{subsec:a5-mnop-two-traces-agree}

\textsc{Attack.} The essay's ``DT and GW as two traces on the same
$\cF_X$'' requires the MNOP change of variable $-q = e^{iu}$, the
\emph{GW/DT correspondence} of MNOP I §§5--6, to make the
identification literal. But MNOP I's statement is a conjecture
(rank-one, at all toric CY$_3$, Conjecture~3R of MNOP I), proved in
stages (MNOP II for smooth projective toric; Pandharipande--Thomas
2009 \texttt{arXiv:0903.3121} stable-pair reformulation;
Toda 2010 \texttt{arXiv:0806.0062} wall-crossing bridge). It is
\emph{not} a priori an ``$\Ethree$-FA trace identity.''

Further: the two traces $\Tr_{\mathrm{cat}}$ and $\Tr_{\mathrm{geom}}$
as defined in A2 and A4 compute the same generating function
$M(-q) = M(e^{iu})|_{-q = e^{iu}}$, so they agree by MNOP \emph{as
functions of the partition variable}; but ``agree as $\Ethree$-FA
traces on $\cF_X$'' is a \emph{stronger} statement about the
underlying trace maps, requiring identification at the level of
the trace map on every element of $\cF_X$, not just on a single
distinguished element.

The bookkeeping: $\Tr_{\mathrm{cat}}$ was computed on $\cO_{\CC^3}$
(the structure sheaf); $\Tr_{\mathrm{geom}}$ was computed on
``holomorphic maps'' (a different object, the mapping-space
insertion). So their agreement is not tautological: it requires a
specific identification of $\cO_X$ with a mapping-space insertion
via the $B$-model equivalence.

\noindent
\textbf{Ghost theorem.} On a toric CY$_3$ $X$, the two $\Ethree$-FA
traces $\Tr^{\cF_X}_{\mathrm{cat}}$ and $\Tr^{\cF_X}_{\mathrm{geom}}$
on $\cF_X = \Obs^q_{\hCS}(X)$ are \emph{equal as functionals on the
localised equivariant cohomology}, under the identification
$-q = e^{iu}$ of DT / GW partition variables. The equality is
literally MNOP for the particular insertions $\cO_X$ (structure
sheaf $\leftrightarrow$ DT side) and $\mathrm{id}_{\cM_g(X)}$
(identity on stable-map moduli $\leftrightarrow$ GW side);
extension to general insertions is the content of MNOP I Conj.\ 3R
in its strongest form.

\noindent
\textbf{Precise error prevented.}
(a) Conflating ``two traces agree on one element'' with
``two traces are equal as maps.'' MNOP as originally stated is the
former (the two partition functions are equal as power series); the
latter is stronger and requires the BCOV $B$-model formulation.
(b) Ignoring the Behrend-sign $(-1)$: MNOP is $(-q)^{\DT} =
(e^{iu})^{\GW}$, with $q\to -q$ reflecting the virtual-dimension sign
$(-1)^{\dim_{\mathrm{vir}}}$. Without it, the partition functions
are off by $\mathbb M(-q)^{\chitop} \neq M(q)^{\chitop}$.
(c) Forgetting the reducedness requirement at $K3\times E$:
MNOP I is for smooth projective toric; the compact non-toric
$K3\times E$ needs MPT 2010 reduction.

\noindent
\textbf{Correct relationship.}

\begin{theorem}[MNOP as agreement of two $\Ethree$-FA traces,
Nekrasov-audited]
\label{wn:thm:mnop-e3-trace-agreement}\ClaimStatusTheorem
Let $X$ be a smooth toric CY$_3$ (e.g.\ $\CC^3$, resolved conifold,
local $\CC\PP^2$, local $\CC\PP^1\times\CC\PP^1$). On
$\cF_X = \Obs^q_{\hCS}(X)$:
\begin{enumerate}[label=\textup{(\roman*)}]
\item $\Tr^{\cF_X}_{\mathrm{cat}}(\cO_X; q) \;=\;
Z^{\DT}_X(q) \;=\; Z^{\DT}_X(-(-q))$ (MacMahon-type function
of $q$, with Behrend sign).
\item $\Tr^{\cF_X}_{\mathrm{geom}}(\mathrm{holomorphic\ maps}; u)
\;=\; Z^{\GW}_X(u) \;=\; \exp\Bigl(\sum_g u^{2g-2} F_g^{\GW}(X)\Bigr)$.
\item Under the change of variable $-q = e^{iu}$ (MNOP I Conj.\ 3R,
proved for smooth toric by MNOP II; PT 2009; Toda 2010),
\[
\Tr^{\cF_X}_{\mathrm{cat}}(\cO_X; q)\bigr|_{-q = e^{iu}}
\;=\; \Tr^{\cF_X}_{\mathrm{geom}}(\mathrm{holomorphic\ maps}; u).
\]
\end{enumerate}
At $X = \CC^3$: both traces equal $M(e^{iu}) = M(-q)|_{-q = e^{iu}}$
explicitly.
\end{theorem}

\begin{proof}[Sketch]
(i): Theorem~\ref{wn:thm:dt-e3-trace-c3-nekrasov}(ii) (for $\CC^3$);
MNOP I Thm.\ 1 (general toric).
(ii): Theorem~\ref{wn:thm:gw-e3-trace-c3-nekrasov}(ii) (for $\CC^3$);
MNOP II Thm.\ 2 (general toric).
(iii): The change of variable is MNOP I Eq.\ (22) / MNOP II Thm.\ 1;
its $\Ethree$-FA-trace formulation is the factorisation-algebra
trace axiom of Costello--Francis--Gwilliam plus the $B$-model duality
of Costello 2011 §6 (tree-level) and Costello 2016 §5 (all-genus).
\end{proof}

\begin{theorem}[MNOP at $K3\times E$, primitive $\beta$: two traces
agree as covariant Siegel modular forms]
\label{wn:thm:mnop-k3e-n1}\ClaimStatusTheorem
For $X = S\times E$ with $\beta\in H_2(S,\ZZ)$ primitive, reduced
classes (MPT 2010), and $(q, t, p) = (e^{2\pi i\tau}, e^{2\pi iz},
e^{2\pi i\sigma})$:
\begin{enumerate}[label=\textup{(\roman*)}]
\item $\Tr^{\cF_X}_{\mathrm{cat}, \mathrm{red}, \mathrm{prim}}(\cO_X;
q, t, p) = -C/\Delta_5(\tau, z, \sigma)^2$.
\item $\Tr^{\cF_X}_{\mathrm{geom}, \mathrm{red}, \mathrm{prim}}(
\mathrm{holomorphic\ maps}; u, \tau, z) =
\exp(\sum_g u^{2g-2} F_g^{\GW, \mathrm{red}, \mathrm{prim}}(S\times
E))$, the reduced primitive GW partition function (MPT 2010 Thm.\ 5).
\item Under $(-q, u) = (e^{iu}, u)$ change of variable (KKV
\texttt{arXiv:hep-th/9910182} plus MPT 2010 Conj.\ 2), the two traces
agree. The resulting function is the Igusa cusp form lift
$-C/\Delta_5^2$ on $\mathbb H_2/\mathrm{Sp}_4(\ZZ)$ (OP 2018 Thm.\ 1).
\end{enumerate}
\end{theorem}

\emph{Three independent verifications at $X=\CC^3$.}
(i)~Plane partitions: $\sum_n |\mathrm{PP}(n)|\,q^n = M(q)$ (MacMahon
1916). (ii)~Topological vertex: $C_{\emptyset\emptyset\emptyset}(q)
= M(q)$ (AKMV 2005). (iii)~GW/DT: MNOP II Thm.\ 2 via
$-q = e^{iu}$. Three routes, one $M(e^{iu})$.

\emph{Three independent verifications at $X=S\times E$, $N=1$.}
(i)~Noether--Lefschetz reduced-class DT: OP 2018 Thm.\ 1.
(ii)~KKV multiple-cover reduction (KKV 1999, as reformulated in
Pandharipande--Thomas 2009): $F^{\GW,\mathrm{red},\mathrm{prim}}_g
(S\times E)$ in terms of K3 elliptic genus $\phi_{0,1}$.
(iii)~Borcherds lift of $\phi_{0,1}$: Lorgat 2020 Thm.\ 4 gives
$\Delta_5 = \mathrm{Bo}(\phi_{0,1})$. Three routes, one
$-C/\Delta_5^2$.

\textsc{Heal.} In the essay, state MNOP as
Theorem~\ref{wn:thm:mnop-e3-trace-agreement} (toric case) and
Theorem~\ref{wn:thm:mnop-k3e-n1} ($K3\times E$, primitive, $N=1$)
side by side. Name the two traces $\Tr_{\mathrm{cat}}$ and
$\Tr_{\mathrm{geom}}$ on the common $\cF_X = \Obs^q_{\hCS}(X)$.
Cite MNOP I, MNOP II, MPT 2010, OP 2018, KKV 1999, Costello 2011
§6, Costello 2016 §5, Costello--Francis--Gwilliam
\texttt{arXiv:1612.00566} §5. Flag at $K3\times E$, $N\ge 2$ as
conjectural (Conj.~1 of Lorgat 2020, conditional on OP 2014
Conj.\ B).

%--------------------------------------------------------------------------
\subsection*{Scope guard: propagation of the MacMahon / $M(-q) = M(e^{iu})$
identity}
\label{subsec:scope-guard-mnop}

\emph{Universal propagation warning.}
The MacMahon identity $Z^{\DT}_X = M(-q)^{\chitop(X)}$ holds
in the following settings, organised by target geometry:
\begin{itemize}
\item $\CC^3$: $T^3$-toric with a single fixed point. MNOP I
Thm.\ 1 directly; $M(-q)^1$.
\item Resolved conifold: $T^3$-toric with compact $\CC\PP^1$.
MNOP I §5 plus Szendrői 2008 \texttt{arXiv:0705.3419}
noncommutative DT; $M(-q)^2$ (signed-Euler 2 from the two
vertices) times the wall-crossing $\prod(1-(-q)^nQ)^n
(1-(-q)^nQ^{-1})^n$ factor. Non-perturbatively the conifold has
extra wall-crossing data (Kontsevich--Soibelman pentagon in
$U_q(\mathfrak{sl}_2)$); the $\Omega$-background perturbative
universality does \emph{not} capture this.
\item Local $\CC\PP^2$, local $\CC\PP^1\times\CC\PP^1$: $T^3$-toric,
compact divisor(s). $M(-q)^{\chitop}$ with $\chitop = 3$
(resp.\ $4$) times wall-crossing factors; genus $F_g$ computable
via topological vertex.
\end{itemize}

\emph{Compact non-toric CY$_3$: MacMahon does not apply.}
\begin{itemize}
\item $K3\times E$: no $T^3$-action. DT trace via reduced virtual
class (MPT 2010) + Noether--Lefschetz stability. Result is
$-C/\Delta_5^2$, paramodular Siegel modular object, \emph{not}
$M(-q)^{\mathrm{anything}}$.
\item Quintic, bicubic, complete-intersection CY$_3$: no
$T^3$-toric structure. DT trace via Behrend-weighted Euler of
Hilbert schemes; generating function not a MacMahon power.
BPS counts via Gopakumar--Vafa $n^g_\beta$, checked for low $\beta$
by direct enumeration (Klemm--Pandharipande 2008
\texttt{arXiv:0705.1653}).
\end{itemize}

\emph{The $\Omega$-background is a toric phenomenon.}
Nekrasov's $(\epsilon_1, \epsilon_2)$-deformation is defined on
$T^2$-equivariant moduli of instantons on $\CC^2$, lifted to
$T^3$-equivariant moduli on toric CY$_3$ via geometric
engineering. It does \emph{not} extend to compact $K3\times E$ or
the quintic; attempts to do so (e.g.\ via ``effective
$\Omega$-background'' on K3-covers) recover only the primitive-class
KKV reduction, not a universal integration tool. This is
CLAUDE.md cache-entry \#1 (toric $\to$ all CY) and is enforced
at every theorem above.

%--------------------------------------------------------------------------
\subsection*{Healed statement: DT and GW as two $\Ethree$-FA traces on
$\cF_X = \PhiFA_3(\cC)$}
\label{subsec:healed-dt-gw-nekrasov}

\begin{theorem}[DT as $\Ethree$-FA categorical trace on
$\cF_X = \PhiFA_3(\Perf(X))$, Nekrasov-audited]
\label{wn:thm:dt-e3-trace-final}\ClaimStatusTheorem
Let $X$ be a smooth CY$_3$; let $\cC = \Perf(X)$; let $\cF_X =
\PhiFA_3(\cC) = \Obs^q_{\hCS}(X)$ the $\Ethree$-factorisation
algebra of observables of 6d holomorphic Chern--Simons on $X$
(Theorem~\ref{wn:thm:plat-hCS-quantum}). The \emph{categorical trace}
$\Tr^{\cF_X}_{\mathrm{cat}}: \cF_X \to K[[q, Q]]$ is the
Costello--Francis--Gwilliam $\Ethree$-observable trace (CFG~§5)
evaluated on the structure-sheaf insertion. Then:
\begin{enumerate}[label=\textup{(\alph*)}]
\item ($\CC^3$, toric, theorem) $\Tr^{\cF_{\CC^3}}_{\mathrm{cat}}
(\cO_{\CC^3}; q) = M(-q)$, the signed MacMahon function
(MNOP I Thm.\ 1; cf.\ Theorem~\ref{wn:thm:dt-e3-trace-c3-nekrasov}).
Three verifications: plane partitions, topological vertex,
$T^3$-equivariant localisation on $\Hilb^n(\CC^3)^T$.
\item (toric $X$, theorem) $\Tr^{\cF_X}_{\mathrm{cat}}(\cO_X; q, Q)
= M(-q)^{\chitop(X)} \cdot W_X(q, Q)$ with $W_X$ the
wall-crossing factor tracking compact divisor classes; MNOP I
Thm.\ 1 plus Szendrői 2008 / Pandharipande--Thomas 2009 at the
conifold, local $\CC\PP^2$, local $\CC\PP^1\times\CC\PP^1$.
\item ($K3\times E$, primitive $\beta$, $N=1$, theorem) on the
reduced class $\Tr^{\cF_{K3\times E}}_{\mathrm{cat}, \mathrm{red},
\mathrm{prim}}(\cO_{K3\times E}; \tau, z, \sigma) = -C/\Delta_5(\tau,
z, \sigma)^2$ (OP 2018 Thm.\ 1; Theorem~\ref{wn:thm:k3e-dt-e3-trace-n1}).
\item ($K3\times E$, primitive $\beta$, $N\in\{2, 3, 4, 6\}$,
conjectural) on the reduced twisted class $\Tr^{\cF_{(K3\times E)/
\langle g_N\rangle}}_{\mathrm{cat}, \mathrm{red}, \mathrm{prim}}(
\cO_{(K3\times E)/\langle g_N\rangle}; \tau, z, \sigma) =
-C_N/F_N(\tau, z, \sigma)^2$ (Lorgat 2020 Conj.\ 1; Wave 13 A5
Table~A2). Status: BOP--Yin 2018 Thm.\ 7 conditional on
OP 2014 Conj.\ B.
\end{enumerate}
\end{theorem}

\begin{theorem}[GW as $\Ethree$-FA geometric trace on $\cF_X =
\PhiFA_3(\Perf(X))$, Nekrasov-audited]
\label{wn:thm:gw-e3-trace-final}\ClaimStatusTheorem
With $X, \cC, \cF_X$ as above, define the \emph{geometric trace}
$\Tr^{\cF_X}_{\mathrm{geom}}(\cdot; u): \cF_X \to K[[u]]$ as the
CFG $\Ethree$-observable trace composed with the holomorphic-map
insertion (= identity on stable-map moduli $\cM_{g, 0}(X,
\beta)$, all $g, \beta$, topological-string coupling $u$). Then:
\begin{enumerate}[label=\textup{(\alph*)}]
\item ($\CC^3$, theorem) $\Tr^{\cF_{\CC^3}}_{\mathrm{geom}}(
\mathrm{holomorphic\ maps}; u) = M(e^{iu})$ (MNOP II Thm.\ 2;
AKMV 2005; Theorem~\ref{wn:thm:gw-e3-trace-c3-nekrasov}).
\item (toric $X$, theorem) $\Tr^{\cF_X}_{\mathrm{geom}}(
\mathrm{holomorphic\ maps}; u, t) = \exp(\sum_g u^{2g-2} F_g^{\GW}
(X, t))$, the all-genus GW free-energy exponential (MNOP II;
BCOV 1994).
\item ($K3\times E$, primitive $\beta$, $N=1$, theorem on reduced
class) identical under $(-q, e^{iu})$ variable change to the
categorical trace; the result is $-C/\Delta_5^2$
(MPT 2010 Thm.\ 5; KKV 1999 reduction).
\end{enumerate}
\end{theorem}

\begin{theorem}[MNOP = trace agreement on $\cF_X = \PhiFA_3(\cC)$]
\label{wn:thm:mnop-trace-equality-nekrasov}\ClaimStatusTheorem
For $X$ a toric CY$_3$, the two $\Ethree$-FA traces
$\Tr^{\cF_X}_{\mathrm{cat}}$ and $\Tr^{\cF_X}_{\mathrm{geom}}$ on
$\cF_X = \PhiFA_3(\Perf(X))$ agree on the distinguished pair
$(\cO_X, \mathrm{holomorphic\ maps})$ under the change of variable
$-q = e^{iu}$:
\[
\Tr^{\cF_X}_{\mathrm{cat}}(\cO_X; q)\bigr|_{-q = e^{iu}}
\;=\; \Tr^{\cF_X}_{\mathrm{geom}}(\mathrm{holomorphic\ maps}; u).
\]
This is MNOP I Conj.\ 3R = MNOP II Thm.\ 1 in the smooth toric case.
At $X = S\times E$ with primitive $\beta$ and reduced classes, the
analogous statement holds (MPT 2010 Thm.\ 5; OP 2018 Thm.\ 1); at
$X = (S\times E)/\langle g_N\rangle$ with $N\in\{2, 3, 4, 6\}$, the
statement is Conjecture~1 of Lorgat 2020.
\end{theorem}

\begin{remark}[Scope and voice]
\label{rem:scope-voice-nekrasov}
The essay's identification
\emph{``DT = categorical trace, GW = geometric trace, on the same
$\cF_X$''} is rigorous as stated in
Theorems~\ref{wn:thm:dt-e3-trace-final}--\ref{wn:thm:mnop-trace-equality-nekrasov}
on the following settings:
\begin{enumerate}
\item Toric $X\in\{\CC^3, \mathrm{conifold}, \mathrm{local}\,\CC\PP^2,
\mathrm{local}\,\CC\PP^1\times\CC\PP^1\}$: theorems (a) and (b)
of both are established by MNOP I, II plus Szendrői 2008 /
PT 2009.
\item $K3\times E$, primitive $\beta$, $N=1$, reduced: theorem (c)
of both is established by OP 2018 plus MPT 2010.
\item $K3\times E$, primitive $\beta$, $N\in\{2, 3, 4, 6\}$, reduced:
conjectural (Lorgat 2020 Conj.\ 1 / BOP--Yin 2018 / OP 2014
Conj.\ B).
\end{enumerate}
Beyond these scopes --- imprimitive $\beta$, $N\in\{5, 7\}$ (Nikulin
obstruction; Wave 13 A5 §``Scope guard''), quintic and bicubic and
general compact non-toric CY$_3$ --- the two-trace identification is
an open programme, not a theorem. The $\Omega$-background
universality of Nekrasov 2002 is perturbatively true on toric
geometries but does not extend to the compact non-toric setting.
Nekrasov's voice throughout: equivariant integration, residues,
Behrend-weighted signs, scope-guards. No $\Omega$-background shortcuts
at $K3\times E$; no MacMahon claims outside of toric scope; no
conflation of the three ``$\Omg(\gamma)$'' invariants.
\end{remark}

%--------------------------------------------------------------------------
\subsection*{Summary: five attack$\leftrightarrow$heal pairs}

\begin{enumerate}
\item[(A1$\mid$H1)] ``Counting in $\cC$'' via $\Omg^{\KS}$ (BPS count,
$\QQ$-valued), coinciding with $\chi^B$ (Behrend-weighted) on
torsion-free classes (Davison 2017). The sum is over numerical
classes in a specified Bridgeland stability chamber; wall-crossing
is Kontsevich--Soibelman pentagon.
\item[(A2$\mid$H2)] MacMahon at $\CC^3$ is the $\Ethree$-FA
categorical trace on $\cF_{\CC^3} = \Obs^q_{\hCS}(\CC^3, \fgl_1)$;
$\chitop(\CC^3) = 1$ in $T^3$-equivariant localisation.
$Z^{\DT}_{\CC^3}(q) = M(-q)$; three independent paths (plane
partitions / localisation / topological vertex).
\item[(A3$\mid$H3)] $Z^{\DT}(K3\times E)|_{\mathrm{red}, \mathrm
{prim}} = -C/\Delta_5^2$ is the global trace on the Igusa quotient;
the fibre trace is the K3 elliptic genus $\phi_{0,1}$; the Borcherds
lift connects them. Scope: primitive $\beta$, reduced class, $N=1$.
Theorem. Rows $N\in\{2, 3, 4, 6\}$ conjectural.
\item[(A4$\mid$H4)] GW as $\Ethree$-FA geometric trace:
Bochner--Martinelli propagator on the hCS path integral,
topological-vertex resummation, all-genus free-energy exponential.
At $\CC^3$: $Z^{\GW}_{\CC^3}(u) = M(e^{iu})$.
\item[(A5$\mid$H5)] MNOP as trace agreement:
$\Tr^{\cF_X}_{\mathrm{cat}}(\cO_X)|_{-q = e^{iu}} =
\Tr^{\cF_X}_{\mathrm{geom}}(\mathrm{holomorphic\ maps})(u)$. Theorem
for smooth toric (MNOP I--II, PT, Toda) and for $K3\times E$
primitive-reduced $N=1$ (OP 2018 + MPT 2010). Conjectural for
$K3\times E$ $N\ge 2$.
\end{enumerate}

Three parallel columns crystallise the essay's
``same $\cF_X$, two traces'' structure:

\begin{center}
\renewcommand{\arraystretch}{1.25}
\begin{tabular}{llll}
\toprule
Target & $\cF_X = \PhiFA_3(\Perf(X))$
& $\Tr^{\cF_X}_{\mathrm{cat}}(\cO_X)$
& $\Tr^{\cF_X}_{\mathrm{geom}}(\mathrm{hol.maps})$\\
\midrule
$\CC^3$ & $\Obs^q_{\hCS}(\CC^3, \fgl_1)$ & $M(-q)$ & $M(e^{iu})$ \\
Conifold & $\Obs^q_{\hCS}(X_{\mathrm{con}}, \fgl_1)$
  & $M(-q)^2 W_{\mathrm{con}}$ & $\exp(\sum_g u^{2g-2} F_g^{\mathrm{con}})$ \\
$K3\times E$, $N=1$ & $\Obs^q_{\hCS}(S\times E, \fgl_1)_{\mathrm{red}}$
  & $-C/\Delta_5^2$ & $-C/\Delta_5^2$ \\
$K3\times E$, $N\in\{2, 3, 4, 6\}$ & twisted, reduced
  & $-C_N/F_N^2$ (conj.) & $-C_N/F_N^2$ (conj.) \\
\bottomrule
\end{tabular}
\end{center}

Row 1 is MNOP I--II theorem; row 2 is MNOP + Szendrői; row 3 is
OP 2018 + MPT 2010 theorem; row 4 is Lorgat 2020 Conjecture~1.
Scope: $N\in\{1, 2, 3, 4, 6\}$ via Nikulin 1980 (AP-CY scope); the
$\Omega$-background apparatus, Nekrasov partition function, and
MacMahon generating function all propagate within the toric column
(rows 1, 2) and \emph{not} into the compact non-toric column
(rows 3, 4), where reduced virtual classes (MPT 2010) and
Noether--Lefschetz stability replace $T^3$-localisation.
The two $\Ethree$-FA traces agree under $-q = e^{iu}$ on the
distinguished pair (structure sheaf $\leftrightarrow$
holomorphic maps), which is MNOP; the underlying $\Ethree$-FA
$\cF_X = \PhiFA_3(\cC)$ is common to both sides, which is the
essay's claim, here made precise.

%--------------------------------------------------------------------------
\subsection*{References}

\begin{itemize}
\item D.~Maulik, N.~Nekrasov, A.~Okounkov, R.~Pandharipande,
``Gromov--Witten theory and Donaldson--Thomas theory'' I, II
(\emph{arXiv:math/0312059}, \emph{arXiv:math/0406092}, 2006).
\item D.~Maulik, R.~Pandharipande, R.~Thomas,
``Curves on $K3$ surfaces and modular forms''
(\emph{arXiv:1007.3085}, 2010).
\item G.~Oberdieck, A.~Pixton, ``Holomorphic anomaly equations
and the Igusa cusp form conjecture'' (\emph{arXiv:1706.10100}, 2018).
\item G.~Oberdieck, R.~Pandharipande, ``Curve counting on
$K3\times E$'' (\emph{arXiv:1411.1514}, 2014).
\item J.~Bryan, G.~Oberdieck, R.~Pandharipande, S.-T.~Yau,
``Donaldson--Thomas theory of $\CC^2\times K3$''
(\emph{arXiv:1802.09180}, 2018).
\item M.~Kontsevich, Y.~Soibelman, ``Stability structures, motivic
Donaldson--Thomas invariants and cluster transformations''
(\emph{arXiv:0811.2435}, 2008).
\item M.~Kontsevich, Y.~Soibelman, ``Cohomological Hall algebra,
exponential Hodge structures and motivic Donaldson--Thomas
invariants'' (\emph{arXiv:1006.2706}, 2010).
\item K.~Behrend, ``Donaldson--Thomas invariants via microlocal
geometry'' (\emph{arXiv:math/0507523}, 2005).
\item B.~Davison, ``Positivity of motivic refined DT invariants''
(\emph{arXiv:1711.04948}, 2017).
\item D.~Joyce, Y.~Song, ``A theory of generalized Donaldson--Thomas
invariants'' (\emph{arXiv:0810.5645}, 2008).
\item R.~Pandharipande, R.~Thomas, ``Curve counting via stable pairs
in the derived category'' (\emph{arXiv:0707.2348}, 2007; PT
\emph{arXiv:0903.3121}, 2009).
\item M.~Aganagic, A.~Klemm, M.~Mariño, C.~Vafa, ``The topological
vertex'' (\emph{arXiv:hep-th/0305132}, 2005).
\item A.~Okounkov, N.~Reshetikhin, ``Random skew plane partitions
and the Pearcey process'' and ``Correlation function of Schur
process with application to local geometry of a random 3-dimensional
Young diagram'' (\emph{arXiv:hep-th/0309208}, 2003).
\item N.~Nekrasov, ``Seiberg--Witten prepotential from instanton
counting'' (\emph{arXiv:hep-th/0206161}, 2002).
\item N.~Nekrasov, A.~Okounkov, ``Seiberg--Witten theory and random
partitions'' (\emph{arXiv:math/0601684}, 2006).
\item K.~Costello, ``Notes on supersymmetric and holomorphic field
theories in dimensions 2 and 4''; ``Holography and Koszul duality:
the example of the M2 brane''
(\emph{arXiv:1303.2632}, 2013; \emph{arXiv:1610.04144}, 2016).
\item K.~Costello, J.~Francis, O.~Gwilliam,
``Factorization algebras in quantum field theory'' Vols.\ I--II
(CUP 2017, 2021).
\item S.~Gukov, A.~Iqbal, C.~Kozcaz, C.~Vafa, ``Link homologies
and the refined topological vertex''
(\emph{arXiv:0705.1368}, 2007).
\item M.~Bershadsky, S.~Cecotti, H.~Ooguri, C.~Vafa,
``Kodaira--Spencer theory of gravity and exact results for quantum
string amplitudes'' (\emph{arXiv:hep-th/9309140}, 1994).
\item S.~Katz, A.~Klemm, C.~Vafa, ``Geometric engineering of
quantum field theories'' (\emph{arXiv:hep-th/9609239}, 1996);
``M-theory, topological strings and spinning black holes''
(\emph{arXiv:hep-th/9910181}, 1999) [KKV].
\item B.~Szendrői, ``Non-commutative Donaldson--Thomas invariants
and the conifold'' (\emph{arXiv:0705.3419}, 2008).
\item R.~Lorgat, ``BKM superalgebras from Borcherds lifts of
twined K3 elliptic genera'' (April 2020, Lorgat paper).
\item Wave 12 A4/A5, Wave 13 A5, Wave 14 companions in
\texttt{notes/}.
\item Theorem~\ref{wn:thm:plat-two-stage} (two-stage factorisation of
$\Phi_3$); Theorem~\ref{wn:thm:plat-hCS-quantum} (quantum
$\Ethree$-structure on $\Obs^q_{\hCS}$).
\end{itemize}
